%% P16 --- Jakobian, reguła łańcuchowa i różniczkowanie automatyczne
%%
%% Zalążek, wygenerowany przez tools/gen_stubs.py z tools/programs.json.
%% Poniższy opis jest kontraktem tego programu: co ma uzasadnić i co ma dać
%% czytelnikowi. Napisanie programu oznacza usunięcie bloku \programstub{} i
%% zastąpienie go programem -- po czym ten plik nie jest już generowany.

\program{Jakobian, reguła łańcuchowa i różniczkowanie automatyczne}
\label{prog:P16}

\programstub{%
\textit{[Opis do przetłumaczenia --- patrz wydanie angielskie.]}

Argument: for vector-valued functions the chain rule multiplies Jacobians;
nobody forms a Jacobian; forward mode computes a Jacobian--vector product
and reverse mode a vector--Jacobian product, and which is cheaper depends
only on the shape of the problem. Payoff: **what \code{loss.backward()}
actually computes**, in full: why a scalar loss with many parameters makes
reverse mode cheaper by a factor of the parameter count; why reverse mode
must keep the activations, so memory scales with depth and batch; gradient
checkpointing as an explicit trade of recomputation for memory, with the
arithmetic done; why a numerical finite-difference gradient is a testing
tool and not an implementation; the three places autodiff silently gives you
something other than the derivative you meant (a non-differentiable point, a
detached tensor, an in-place write).

\medskip
\textbf{Planowana długość:} 65 ramek.
}
